\section*{第 \ref{ch:b1:units-dimensions} 章——单位、量纲与测量}

\begin{solution}{exo:b1:units-dimensions:1}
$\unit{J} = \unit{kg.m^2.s^{-2}}$; $\unit{W} = \unit{kg.m^2.s^{-3}}$;
$\unit{Pa} = \unit{kg.m^{-1}.s^{-2}}$; $\unit{C} = \unit{A.s}$.
$\unit{Pa.m^3} = \unit{kg.m^{-1}.s^{-2}.m^3} = \unit{kg.m^2.s^{-2}} = \unit{J}$
（压力做的功，$P\,\dd V$）。
\end{solution}

\begin{solution}{exo:b1:units-dimensions:2}
压强 $M L^{-1} T^{-2}$；功率 $M L^2 T^{-3}$；电荷 $I\,T$；
电压 $= $ 功率/电流 $= M L^2 T^{-3} I^{-1}$。
$[G] = [F r^2/m^2] = M L T^{-2} L^2 M^{-2} = L^3 M^{-1} T^{-2}$;
$[h] = [E/\nu] = M L^2 T^{-2} \cdot T = M L^2 T^{-1}$;
$[k_B] = [E/T] = M L^2 T^{-2} \Theta^{-1}$.
\end{solution}

\begin{solution}{exo:b1:units-dimensions:3}
$\sqrt{2gh}$：$\sqrt{L T^{-2} \cdot L} = L T^{-1}$，齐次。
$\tfrac12 mv$：$M L T^{-1} \neq M L^2 T^{-2}$，不是能量。
$\rho g h$：$M L^{-3} \cdot L T^{-2} \cdot L = M L^{-1} T^{-2}$，是压强。
$v_0 t + \tfrac12 g t^2$：两项都是 $L$，齐次。
$2\pi\sqrt{g/\ell}$：$\sqrt{T^{-2}} = T^{-1}$——是频率，不是周期。
$I_0\eu^{-t/RC}$：$RC$ 在 $\unit{\ohm.F} = \unit{s}$ 中，指数部分无量纲，齐次。
\end{solution}

\begin{solution}{exo:b1:units-dimensions:4}
直尺：$\delta = \qty{1}{mm}$，$u = 1/\sqrt{12} = \qty{0.29}{mm}$。
天平：$u = 0.01/\sqrt{12} = \qty{0.003}{g}$。
电压表：$a = 0.005 \times 4.80 = \qty{0.024}{V}$，
$u = 0.024/\sqrt3 = \qty{0.014}{V}$.
\end{solution}

\begin{solution}{exo:b1:units-dimensions:5}
$v = k F^\alpha \mu^\beta$: $L T^{-1} = (M L T^{-2})^\alpha (M L^{-1})^\beta$,
于是 $\alpha + \beta = 0$，$\alpha - \beta = 1$，$-2\alpha = -1$：
$\alpha = \tfrac12$，$\beta = -\tfrac12$，$v = k\sqrt{F/\mu}$（实际上 $k = 1$）。
声音：$P/\rho$ 的量纲为 $M L^{-1} T^{-2} / (M L^{-3}) = L^2 T^{-2}$，
故 $v = k\sqrt{P/\rho}$（对空气，$k = \sqrt\gamma \approx 1.2$）。
\end{solution}

\begin{solution}{exo:b1:units-dimensions:6}
总和 $= \qty{3620}{ms}$，$\bar t = \qty{452.5}{ms}$。偏差为
$-0.5, -5.5, 8.5, 2.5, -3.5, 5.5, -8.5, 1.5$；平方和为 $226$；
$s = \sqrt{226/7} = \qty{5.7}{ms}$; $u(\bar t) = 5.7/\sqrt8 = \qty{2.0}{ms}$.
结果：$t = \qty{452.5\pm2.0}{ms}$。
\end{solution}

\begin{solution}{exo:b1:units-dimensions:7}
$R = 4.80/0.0124 = \qty{387}{\ohm}$。相对不确定度为 $0.03/4.80 =
0.63\%$ 和 $0.2/12.4 = 1.6\%$；合成为 $\sqrt{0.63^2 + 1.6^2} = 1.7\%$，
$u(R) = \qty{7}{\ohm}$：$R = \qty{387\pm7}{\ohm}$。电流表的
误差占主导——应首先改进电流表（或其量程）。
\end{solution}

\begin{solution}{exo:b1:units-dimensions:8}
$z = 7/\sqrt{16 + 9} = 1.4 \leq 2$：相容。
第三与第一组比较：$z = 9/\sqrt{4 + 16} = 2.0$：恰在临界处，勉强
相容。第三与第二组比较：$z = 16/\sqrt{4 + 9} = 4.4$：不相容
——两者中至少有一个存在系统误差，或
不确定度被低估。
\end{solution}

\begin{solution}{exo:b1:units-dimensions:9}
$V = \pi d^3/6 = \pi \times 25.40^3/6 = \qty{8580}{mm^3}$.
$u(V)/V = 3\,u(d)/d = 3 \times 0.02/25.40 = 0.24\%$,
$u(V) = \qty{20}{mm^3}$: $V = \qty{8.58\pm0.02}{cm^3}$.
\end{solution}

\begin{solution}{exo:b1:units-dimensions:10}
$F = k\eta^a R^b v^c$: $M$: $a = 1$; $T$: $-a - c = -2$, $c = 1$;
$L$：$-a + b + c = 1$，$b = 1$：$F = k\eta R v$（斯托克斯，$k = 6\pi$）。
若改用 $\rho$：$M$：$a = 1$；$T$：$-c = -2$，$c = 2$；$L$：
$-3 + b + 2 = 1$, $b = 2$: $F = k\rho R^2 v^2$.
雨滴：$\eta R v = 1.8\times10^{-5} \times 10^{-3} \times 5 \approx
\qty{1e-7}{N}$; $\rho R^2 v^2 = 1.2 \times 10^{-6} \times 25 =
\qty{3e-5}{N}$，大了约三百倍：适用惯性（高速）
区——比值 $\rho v R/\eta \approx 330$ 即为
雷诺数。
\end{solution}

\begin{solution}{exo:b1:units-dimensions:11}
$\bar I = \qty{5.0}{mA}$，$\bar U = \qty{1.000}{V}$；$I$ 的偏差为
$-3, -1, 1, 3$；$U$ 的偏差为 $-0.59, -0.22, 0.22, 0.59$。
$\sum (I_i - \bar I)(U_i - \bar U) = 3.98$, $\sum (I_i - \bar I)^2 = 20$:
$a = \qty{0.199}{V/mA} = \qty{199}{\ohm}$, $b = 1.000 - 0.199 \times 5 =
\qty{0.005}{V}$。取每点 $u(U) = \qty{0.02}{V}$，截距
\qty{0.005}{V} 完全在不确定度范围内：与零相容，电阻器
是欧姆元件，$R \approx \qty{199}{\ohm}$。
\end{solution}

\begin{solution}{exo:b1:units-dimensions:12}
$\ell = G^a h^b c^c$: $M$: $-a + b = 0$; $T$: $-2a - b - c = 0$; $L$:
$3a + 2b + c = 1$。于是 $b = a$，$c = -3a$，$2a = 1$：
$\ell_P = \sqrt{Gh/c^3} = \sqrt{4.42\times10^{-44}/2.7\times10^{25}}
= \qty{4.1e-35}{m}$; $t_P = \ell_P/c = \qty{1.4e-43}{s}$;
$m_P = \sqrt{hc/G} = \sqrt{1.99\times10^{-25}/6.67\times10^{-11}}
= \qty{5.5e-8}{kg}$。普朗克长度比
质子小约二十个数量级；普朗克质量 \qty{55}{\micro g}，只是一粒尘埃
——对粒子很大，对人却很小（轻了约 $10^9$ 倍）。
\end{solution}

\begin{solution}{pb:b1:units-dimensions:1}
\textbf{1.} $\sqrt{L/(L T^{-2})} = \sqrt{T^2} = T$.

\textbf{2.} $\tau = k\ell^\alpha m^\beta g^\gamma$ 给出
$T = L^{\alpha+\gamma} M^\beta T^{-2\gamma}$，故 $\beta = 0$：
$\ell$ 与 $g$ 的任何组合都无法抵消质量。

\textbf{3.} $\theta_0$ 无量纲，因此任何因子 $f(\theta_0)$ 都
不影响量纲分析。取较小的振幅（约
$10^\circ$ 以下），此时 $f \approx 1$，误差优于 $0.2\%$。

\textbf{4.} $g = 4\pi^2\ell/\tau^2$。

\textbf{5.} $\ell = 99.0 + 1.00 = \qty{100.0}{cm} = \qty{1.000}{m}$。

\textbf{6.} $u = \qty{1}{mm}/\sqrt{12} = \qty{0.29}{mm}$。

\textbf{7.} 对准误差：$5/\sqrt3 = \qty{2.9}{mm}$；合成
$\sqrt{0.29^2 + 2.9^2} = \qty{2.9}{mm}$——卷尺读数
的影响可以忽略。$u(\ell)/\ell = 0.29\%$。

\textbf{8.} 反应时间误差（约 \qty{0.1}{s}）
无论测量一个还是十个周期，其误差相同，因此对单个周期的影响
小十倍。重复测量可以显示离散程度（由此得到 $s$）
并使平均值的不确定度除以 $\sqrt8$。

\textbf{9.} 总和 $= \qty{160.80}{s}$，$\overline{t_{10}} = \qty{20.100}{s}$。
偏差为 $0.02, -0.05, 0.08, -0.01, 0.05, -0.08, 0.01, -0.02$；平方和
为 $\num{188e-4}$；$s = \sqrt{188\times10^{-4}/7} = \qty{0.052}{s}$。

\textbf{10.} $u(\overline{t_{10}}) = 0.052/\sqrt8 = \qty{0.018}{s}$；
$\tau = \qty{2.0100}{s}$，$u(\tau) = \qty{0.0018}{s}$。

\textbf{11.} $u(\tau)/\tau = 0.09\%$，比
$u(\ell)/\ell = 0.29\%$ 小三倍。

\textbf{12.} $g = 4\pi^2 \times 1.000/2.0100^2 = 39.48/4.040 =
\qty{9.772}{m/s^2}$。

\textbf{13.} $g \propto \ell\,\tau^{-2}$：
$\dfrac{u(g)}{g} = \sqrt{\left(\dfrac{u(\ell)}{\ell}\right)^2
+ \left(2\dfrac{u(\tau)}{\tau}\right)^2} = \sqrt{0.29^2 + 0.18^2}\,\%
= 0.34\%$。

\textbf{14.} $u(g) = 0.0034 \times 9.772 = \qty{0.033}{m/s^2}$：
$g = \qty{9.77\pm0.03}{m/s^2}$。

\textbf{15.} $z = \abs{9.772 - 9.81}/0.033 = 1.2 \leq 2$：相容。

\textbf{16.} 长度项是主要误差来源（$0.29\%$，周期项为 $0.18\%$）。
将两端定位到 $\pm\qty{1}{mm}$ 可使 $u(\ell)/\ell$ 降至
$0.06\%$，$u(g)/g$ 降至 $\sqrt{0.06^2 + 0.18^2} = 0.19\%$；将长度加倍，
使 $u(\ell)/\ell$ 减半，并使 $\tau$ 变长，两项都得到改善；
测量二十个周期而非十个，只使周期项减半：
$u(g)/g = \sqrt{0.29^2 + 0.09^2} = 0.30\%$——收益很小，
因为此时长度才是瓶颈。

\textbf{17.} 否：对每次读数都相同的偏移是
\emph{系统误差}（偏差）；误差预算只涵盖随机离散和已声明的容差，
不包括一致的偏移。它可以通过改变方法或仪器来检测，
或者——见第五部分——通过 $\tau^2$--$\ell$ 直线的非零
截距来检测。

\textbf{18.} $\tau^2 = (4\pi^2/g)\,\ell$ 是过原点的直线：
线性与零截距都可以通过目测和最小二乘法检验；
而 $\tau \propto \sqrt\ell$ 则不能。

\textbf{19.} $\bar\ell = \qty{0.800}{m}$；$\overline{\tau^2} = 16.16/5 =
\qty{3.232}{s^2}$。

\textbf{20.} $\ell$ 的偏差：$-0.4, -0.2, 0, 0.2, 0.4$；
$\tau^2$ 的偏差：$-1.612, -0.822, 0.018, 0.808, 1.608$；交叉和 $= 1.614$，
$\sum(\ell_i - \bar\ell)^2 = 0.400$：$a = \qty{4.035}{s^2/m}$，
$b = 3.232 - 4.035 \times 0.800 = \qty{0.004}{s^2}$。

\textbf{21.} $g = 4\pi^2/a = 39.48/4.035 = \qty{9.78}{m/s^2}$。

\textbf{22.} $b \approx 0$，与模型预测一致。若 $b$ 明显非零，
则说明每个长度中都存在系统偏移 $\delta$
（$\tau^2 = a(\ell_{\text{测}} - \delta)$ 给出 $b = -a\delta$），例如
小球中心定位错误——或者是振幅效应。

\textbf{23.} $u(g)/g = u(a)/a = 0.05/4.035 = 1.2\%$，$u(g) =
\qty{0.12}{m/s^2}$：$g = \qty{9.78\pm0.12}{m/s^2}$。

\textbf{24.} $z = \abs{9.772 - 9.784}/\sqrt{0.033^2 + 0.12^2} = 0.1$：
完全相容。单次精心测量的精度更高；
这条直线验证了模型，并可排查长度偏移。

\textbf{25.} $\tau = 2\pi\sqrt{10/9.8} = \qty{6.3}{s}$；若只测量一次摆动，
反应时间约为 \qty{0.2}{s}，则 $u(\tau)/\tau \approx
3\%$，因此 $u(g)/g \approx 6\%$，而 $u(\ell)/\ell$ 降至
$0.03\%$：比第四部分差二十倍。长度带来的收益
在测量单个周期时被浪费。最佳策略：长摆\emph{且}多个周期
反复测量——对 \qty{10}{m} 的摆测十次摆动，
重复八次，可得 $u(\tau)/\tau \approx 0.03\%$，$u(g)/g \approx 0.07\%$。
\end{solution}
